\section*{Abstract}\label{abstract}
\addcontentsline{toc}{section}{Abstract}

\textbf{From verified criticism to executable repair---and the limits of
economic compounding.} An invention-automation machine must do more than
recognize an analogy, select a familiar solver, or reject an inadequate
answer. A validated failure must change an executable construction; that
construction must improve useful work against competent alternatives;
and complete accounts must establish whether it enables the next
construction. We audit the actual AGI ALPHA developer artifacts and
implement an additive corrective-constraint compiler inside the existing
mission and evidence flow. In a fresh, frozen exploratory panel of 36
finite operational models, the repaired version accepts 36/36 schedules
after generating 594 independently checked constraints. A full
conventional linear program also accepts 36/36 and uses less measured
deployment compute: 0.239909 versus 0.123720 CPU-proxy seconds per
attempted task. The earlier correction version accepted 33/36 because
three repairs exceeded the unchanged CPU limit. Read-only audits
reproduce all 288 new native records and validate all 1,220 generated
constraints across both versions. Separately, 128 saved developer
transfer/successor certificates replay successfully but establish no
accepted-output advantage. Current simulated accounts reconcile an
instrument operating loss of \textbf{8.574947 SIM USD}, despite an
earlier positive checkpoint; unrestricted earned funds are zero. Battery
reuse, difficult textile sorting, and small-building energy services
remain unadmitted opportunities, not verified loss-making cohorts.
Historical eight-arm and learning summaries do not establish practical
synergy or progressive usefulness, and their full raw archive was not
recovered. The result is an executed closure of a narrow
criticism-to-construction gap and a consequential negative audit---not a
demonstrated economic transformation, general invention engine, or
completed flywheel.

\textbf{Keywords:} invention automation; structural transfer;
counterexample-guided construction; productive inheritance;
complete-cost accounting; procedural verification; AGI ALPHA.

\textbf{Author:} Vincent Boucher, President, MONTREAL.AI.
\textbf{Status:} local research edition; not externally peer reviewed,
deployed, or remotely published. The ambition is to turn each verified
invention into a means of discovering, constructing, and economically
enabling the next. That ambition is the hypothesis, not the result.

